\subsection{EXTENSION OF THE RING OF OPERATORS OF A HOMOMORPHISM MODULE}

Let A be a commutative ring, B a ring, \(\rho: \mathbf{A} \to \mathbf{B}\) a ring homomorphism and E, F two A-modules; as B is an (A, A)-bimodule (by means of \(\rho\) ) and F can be considered as an (A, A)-bimodule, there are on the Z-module B \(\otimes_{\mathbf{A}}\) F two A-module structures, under which respectively \(a(b \otimes y) = (\rho(a)b) \otimes y\) and \(a(b \otimes y) = b \otimes (ay)\) for \(a \in \mathbf{A}\) , \(b \in \mathbf{B}\) , \(y \in \mathbf{F}\) . We shall denote the two A-modules thus defined by G' and G''; G' is moreover just the A-module \(\rho_*(\rho^*(\mathbf{F}))\) .

This being so, in the definition of the canonical homomorphism of § 4, no. 2, formula (7), we replace B by A, the B-module F by the ring B considered as an A-module by means of \(\rho\) and G by F considered as an (A, A)-bimodule; as A is commutative, we may write the canonical Z-homomorphism obtained as

\[
\mathbf {B} \otimes_ {\mathbf {A}} \operatorname{Hom} _ {\mathbf {A}} (\mathbf {E}, \mathbf {F}) \rightarrow \operatorname{Hom} _ {\mathbf {A}} (\mathbf {E}, \mathbf {G} ^ {\prime \prime}).\tag{15}
\]

On the other hand (no. 1, formula (2)), there is a canonical Z-isomorphism

\[
\mathrm{Hom} _ {\mathbf {A}} (\mathbf {E}, \mathbf {G} ^ {\prime}) = \mathrm{Hom} _ {\mathbf {A}} (\mathbf {E}, \rho_ {*} (\rho^ {*} (\mathbf {F}))) \rightarrow \mathrm{Hom} _ {\mathbf {B}} (\rho^ {*} (\mathbf {E}), \rho^ {*} (\mathbf {F})). \tag {16}
\]

Suppose now that \(\rho(A)\) is contained in the centre of \(B\) , in which case \(\rho\) is also called a central homomorphism \(*\) (or \(\rho\) is said to define an A-algebra structure on \(B\) , cf.~III, §1, no. 3). Then the A-module structures of \(G'\) and \(G''\) are identical and composing the homomorphisms (16) and (15) we thus obtain a canonical Z-homomorphism

\[
\omega : \mathrm{B} \otimes_ {\mathrm{A}} \operatorname{Hom} _ {\mathrm{A}} (\mathrm{E}, \mathrm{F}) \rightarrow \operatorname{Hom} _ {\mathrm{B}} \left(\mathrm{E} _ {(\mathrm{B})}, \mathrm{F} _ {(\mathrm{B})}\right)\tag{17}
\]

which is characterized by the fact that, for all \(u \in \operatorname{Hom}_{\mathbf{A}}(\mathbf{E}, \mathbf{F})\) and all \(b \in B\)

\[
\omega (b \otimes u) = r _ {b} \otimes u,\tag{18}
\]

where \(r_b\) denotes right multiplication by \(b\) in B.

Moreover, the hypothesis that \(\omega\) is a central homomorphism implies that \((bb')\rho(a) = b\rho(a)b'\) for \(b, b'\) in B and \(a \in A\) ; in other words the right B-module structure of \(B_d\) is compatible with its A-module structure; it thus defines on B \(\otimes_A\) \(\mathrm{Hom}_{\mathbf{A}}(\mathbf{E}, \mathbf{F})\) a right B-module structure (§ 3, no. 4) and also on \(\mathrm{F}_{(\mathbf{B})} = \mathbf{B} \otimes_{\mathbf{A}} \mathbf{F}\) , and finally, as the left and right B-module structures on \(\mathrm{F}_{(\mathbf{B})}\) are compatible, a right B-module structure is also obtained on \(\mathrm{Hom}_{\mathbf{B}}(\mathrm{E}_{(\mathbf{B})}, \mathrm{F}_{(\mathbf{B})})\) (§ 1, no. 14). Then it is immediately verified that (17) is a right B-module homomorphism for these structures.

PROPOSITION 7. Let A be a commutative ring, B a ring, \(\rho: \mathbf{A} \to \mathbf{B}\) a central homomorphism and E, F two A-modules.

\begin{enumerate}
\def\labelenumi{(\roman{enumi})}
\item
  If B is a projective (resp. finitely generated projective) A-module, the homomorphism (17) is injective (resp. bijective).
\item
  If \(\mathbf{E}\) is a finitely generated projective A-module, the homomorphism (17) is bijective.
\end{enumerate}

As (16) is bijective, the proposition follows from § 4, no. 2. Proposition 2, applied to the canonical homomorphism (15).

